\DocumentMetadata{pdfversion=1.4}

\documentclass[12pt, a4paper]{article}
\usepackage[a4paper, top=2cm, bottom=2cm, left=2cm, right=2cm]{geometry}

\usepackage{fontspec}
\setmainfont{Source Serif Pro}
\setmonofont{Iosevka Fixed}

\usepackage{polyglossia}
\setdefaultlanguage{polish}

\usepackage{microtype}
\usepackage[]{hyperref}
\hypersetup{
	colorlinks=false,
	pdfborderstyle={/S/U/W 0.5},
	allbordercolors=black
}

\author{Kamil J. Hantszel, Aleksander Kała}
\title{Opis technologiczny projektu}
\date{}

\begin{document}
\maketitle

\section{Temat projektu}
\textbf{Ray tracer} — własny silnik renderujący statyczne obrazy korzystający z algorytmu śledzenia promieni, pisany
w języku C++.

\section{Repozytorium}
\url{https://github.com/Althorion/cpp-ray-tracer.git}

\section{Opis projektu}
Celem projektu jest stworzenie własnego silnika renderującego statyczne obrazy przy wykorzystaniu algorytmu śledzenia
promieni i napisanie go w sposób maksymalizujący wykorzystanie paradygmatu programowania obiektowego.

Program będzie przyjmował z linii komend adres pliku z opisem sceny (we własnym prostym języku opisu), rozdzielczość
generowanego obrazka i ścieżkę do zapisu wygenerowanego obrazka, po czym wygeneruje i zapisze obrazek.

Grafika generowana będzie w formacie PPM. Możliwe kształty zawarte w renderowanej scenie ograniczone są tylko do kul.
Implementacja będzie obejmować różne materiały kul, z różnymi parametrami odbijania światła (metal, szkło, plastik).

Program będzie zawierał testy automatyczne, wykonane głównie w technice testów własności (i ewentualnie dodatkowych
testów jednostkowych). Będzie obsługiwał proste logowanie i obsługę błędów.

\section{Stos technologiczny}
\begin{itemize}
	\item \textbf{C++} z wykorzystaniem standardu C++26 (w zakresie obsługiwanym przez wykorzystywane kompilatory).
	\item \textbf{CMake} — zarządzanie kompilacją.
	\item \textbf{Clang-Tidy} — linter.
	\item \textbf{ClangFormat} — formatowanie kodu.
	\item \textbf{rapidcheck} — testy parametryczne.
	\item \textbf{pcg-cpp} — szybki generator pseudolosowy.
\end{itemize}

\section{Podział prac w zespole}
\begin{itemize}
	\item \textbf {Kamil J. Hantszel}
		\begin{itemize}
			\item Zarządzanie projektem.
			\item Implementacja głównego algorytmu i struktur danych.
			\item Testowanie ww.
			\item Dokumentacja.            
		\end{itemize}
	\item \textbf {Aleksander Kała}
		\begin{itemize}
			\item CMake
			\item Obsługa wejścia-wyjścia.
			\item Logowanie i obsługa błędów, zwłaszcza w I/O.
			\item Testowanie ww.
		\end{itemize}
\end{itemize}

\end{document}

%%% Local Variables:
%%% ispell-local-dictionary: "pl-PL" 
%%% End:
